\documentclass{ufctex}

% Historical v1 metadata API.
\trabalhoacademico{tese}
\ies{Universidade Federal do Ceará}
\iessigla{UFC}
\centro{Centro de Ciências}
\departamento{Departamento de Computação}
\programadoutorado{Programa de Pós-Graduação em Ciência da Computação}
\nomedodoutorado{Doutorado em Ciência da Computação}
\doutorem{Ciência da Computação}
\areadeconcentracaodoutorado{Computação Gráfica}
\autor{Autor Legado Teste}
\titulo{Documento Legado de Validação}
\data{2026}
\local{Fortaleza}
\orientador{Prof. Orientador Legado}
\orientadories{Universidade Federal do Ceará}
\orientadorcentro{Centro de Ciências}
\orientadorfeminino{nao}
\incluirfichacatalografica{nao}

\ufcbibliografia{tests/fixtures/referencias-v2.bib}

\begin{document}

\imprimircapa
\imprimirfolhaderosto
\imprimirsumario

\textual
\section{Introdução}
A API histórica de citação permanece disponível: \citeonline{silva2020}.

\UFCfig{
  \Caption{\label{fig:compat-v1} Objeto pela API legada}
}{
  \rule{6cm}{2cm}
}{
  \Fonte{Elaboração própria.}
  \Nota{Objeto usado apenas para regressão de compatibilidade.}
}
A Figura~\ref{fig:compat-v1} foi produzida pela interface histórica.

\seprojeto{ERRO-PERFIL-PROJETO}{PERFIL-ACADEMICO-CORRETO}

\imprimirreferencias

\imprimirapendices
\apendice{Material complementar produzido pelo autor}
Conteúdo do apêndice legado.

\imprimiranexos
\anexo{Documento externo}
Conteúdo do anexo legado.

\end{document}
